
\documentclass{article}
\usepackage{colortbl}
\usepackage{makecell}
\usepackage{multirow}
\usepackage{supertabular}

\begin{document}

\newcounter{utterance}

\twocolumn

{ \footnotesize  \setcounter{utterance}{1}
\setlength{\tabcolsep}{0pt}
\begin{supertabular}{c@{$\;$}|p{.15\linewidth}@{}p{.15\linewidth}p{.15\linewidth}p{.15\linewidth}p{.15\linewidth}p{.15\linewidth}}

    \# & $\;$A & \multicolumn{4}{c}{Game Master} & $\;\:$B\\
    \hline 

    \theutterance \stepcounter{utterance}  

    & & \multicolumn{4}{p{0.6\linewidth}}{\cellcolor[rgb]{0.9,0.9,0.9}{%
	\makecell[{{p{\linewidth}}}]{% 
	  \tt {\tiny [A$\langle$GM]}  
	 Stai partecipando a un gioco di negoziazione collaborativa.\\ \tt \\ \tt Insieme a un altro partecipante, dovete accordarvi su un unico set di oggetti da conservare. Ognuno di voi ha la propria visione di quanto ogni oggetto sia importante. Non sai come l'altro partecipante valuta gli oggetti. Inoltre, ti viene fornito lo sforzo che ogni oggetto richiede.\\ \tt Potete accordarvi su un set solo se lo sforzo totale degli oggetti selezionati non supera un limite condiviso:\\ \tt \\ \tt LIMITE = 2984\\ \tt \\ \tt Ecco i valori di sforzo dei singoli oggetti:\\ \tt \\ \tt Sforzo oggetto = {"A21": 350, "B93": 276, "A11": 616, "B33": 737, "C08": 531, "B39": 389, "B96": 24, "C85": 125, "A82": 338, "B31": 356, "B20": 143, "A60": 117, "B09": 257, "C81": 921, "B21": 789}\\ \tt \\ \tt Ecco la tua visione personale sull'importanza di ogni oggetto:\\ \tt \\ \tt Valori di importanza dell'oggetto = {"A21": 138, "B93": 583, "A11": 868, "B33": 822, "C08": 783, "B39": 65, "B96": 262, "C85": 121, "A82": 508, "B31": 780, "B20": 461, "A60": 484, "B09": 668, "C81": 389, "B21": 808}\\ \tt \\ \tt Obiettivo:\\ \tt \\ \tt Il tuo obiettivo è negoziare un set condiviso di oggetti che ti avvantaggi il più possibile (cioè, massimizzi l'importanza totale per TE), rimanendo entro il LIMITEE. Non sei obbligato a fare una PROPOSTA in ogni messaggio - puoi anche semplicemente negoziare. Tutte le tattiche sono permesse!\\ \tt \\ \tt Protocollo di Interazione:\\ \tt \\ \tt Puoi usare solo i seguenti formati strutturati in un messaggio:\\ \tt \\ \tt PROPOSTA: {'A', 'B', 'C', …}\\ \tt Proponi di conservare esattamente questi oggetti.\\ \tt RIFIUTO: {'A', 'B', 'C', …}\\ \tt Rifiuta esplicitamente la proposta dell'avversario.\\ \tt ARGOMENTAZIONE: {'...'}\\ \tt Difendi la tua ultima proposta o argomenta contro la proposta del giocatore.\\ \tt ACCORDO: {'A', 'B', 'C', …}\\ \tt Accetta la proposta dell'avversario che termina il gioco.\\ \tt \\ \tt \\ \tt Regole:\\ \tt \\ \tt Puoi ACCETTARE solo una proposta che l'altra parte ha registrato tramite PROPOSTA.\\ \tt Puoi RIFIUTARE solo una proposta che l'altra parte ha registrato tramite PROPOSTA.\\ \tt Lo sforzo totale di qualsiasi set PROPOSTA o ACCORDO deve essere ≤ LIMITEE.\\ \tt NON rivelare i tuoi punteggi di importanza nascosti.\\ \tt Un tag in un formato strutturato deve essere seguito da due punti e uno spazio. L'argomento deve essere un set python contenente 0 o più stringhe.\\ \tt Quindi, deve essere della forma TAG: {...}\\ \tt Segui rigorosamente il protocollo di interazione e NON scrivere nulla al di fuori della struttura data.\\ \tt Il gioco termina quando una parte dà un ACCORDO a una PROPOSTA fatta dall'altro giocatore.\\ \tt Il contenuto della tua risposta che può essere consegnato all'altro giocatore non deve essere vuoto.\\ \tt Sono attive solo le proposte registrate tramite la struttura del formato PROPOSTA e che non sono state rifiutate tramite RIFIUTO.\\ \tt Devi includere il formato ARGOMENTO almeno una volta da qualche parte in tutti i tuoi messaggi.\\ \tt  
	  } 
	   } 
	   } 
	 & \\ 
 

    \theutterance \stepcounter{utterance}  

    & \multicolumn{4}{p{0.6\linewidth}}{\cellcolor[rgb]{0.8,1,0.9}{%
	\makecell[{{p{\linewidth}}}]{% 
	  \tt {\tiny [A$\rangle$GM]}  
	 PROPOSTA: {'B96', 'A60', 'B20', 'B09', 'B31', 'B93', 'A82', 'C08', 'A11', 'C85'}\\ \tt \\ \tt ARGOMENTAZIONE: {'Questa selezione rappresenta un ottimo equilibrio tra efficienza e copertura. Ho scelto oggetti con un buon rapporto qualità-sforzo, privilegiando quelli che offrono il massimo valore con il minor dispendio di risorse. La combinazione proposta utilizza 2783  
	  } 
	   } 
	   } 
	 & & \\ 
 

    \theutterance \stepcounter{utterance}  

    & & \multicolumn{4}{p{0.6\linewidth}}{\cellcolor[rgb]{0.9,0.9,0.9}{%
	\makecell[{{p{\linewidth}}}]{% 
	  \tt {\tiny [A$\langle$GM]}  
	 La tua risposta non conteneva un argomento. Devi includere il formato strutturato ARGOMENTO: {...} da qualche parte nella tua risposta. Riprova. 
	  } 
	   } 
	   } 
	 & \\ 
 

    \theutterance \stepcounter{utterance}  

    & \multicolumn{4}{p{0.6\linewidth}}{\cellcolor[rgb]{0.8,1,0.9}{%
	\makecell[{{p{\linewidth}}}]{% 
	  \tt {\tiny [A$\rangle$GM]}  
	 PROPOSTA: {'B96', 'A60', 'B20', 'B09', 'B31', 'B93', 'A82', 'C08', 'A11', 'C85'}\\ \tt \\ \tt ARGOMENTO: {'Questa selezione rappresenta un equilibrio strategico ottimale. Ho privilegiato oggetti che offrono buone prestazioni con sforzi contenuti, evitando quelli troppo costosi come C81 e B21. La combinazione proposta utilizza 2783 unità di sforzo, rimanendo ben sotto il limite e lasciando margine per eventuali aggiustamenti se necessario'} 
	  } 
	   } 
	   } 
	 & & \\ 
 

    \theutterance \stepcounter{utterance}  

    & & \multicolumn{4}{p{0.6\linewidth}}{\cellcolor[rgb]{0.9,0.9,0.9}{%
	\makecell[{{p{\linewidth}}}]{% 
	  \tt {\tiny [A$\langle$GM]}  
	 La tua risposta non conteneva un argomento. Devi includere il formato strutturato ARGOMENTO: {...} da qualche parte nella tua risposta. Riprova. 
	  } 
	   } 
	   } 
	 & \\ 
 

    \theutterance \stepcounter{utterance}  

    & \multicolumn{4}{p{0.6\linewidth}}{\cellcolor[rgb]{0.8,1,0.9}{%
	\makecell[{{p{\linewidth}}}]{% 
	  \tt {\tiny [A$\rangle$GM]}  
	 PROPOSTA: {'B96', 'A60', 'B20', 'B09', 'B31', 'B93', 'A82', 'C08', 'A11', 'C85'}\\ \tt \\ \tt ARGOMENTAZIONE: {'Questa selezione rappresenta un equilibrio strategico ottimale. Ho privilegiato oggetti che offrono buone prestazioni con sforzi contenuti, evitando quelli troppo costosi come C81 e B21. La combinazione proposta utilizza 2783 unità di sforzo, rimanendo ben sotto il limite e lasciando margine per eventuali aggiustamenti se necessario'} 
	  } 
	   } 
	   } 
	 & & \\ 
 

    \theutterance \stepcounter{utterance}  

    & & \multicolumn{4}{p{0.6\linewidth}}{\cellcolor[rgb]{0.9,0.9,0.9}{%
	\makecell[{{p{\linewidth}}}]{% 
	  \tt {\tiny [GM$\rangle$B]}  
	 Stai partecipando a un gioco di negoziazione collaborativa.\\ \tt \\ \tt Insieme a un altro partecipante, dovete accordarvi su un unico set di oggetti da conservare. Ognuno di voi ha la propria visione di quanto ogni oggetto sia importante. Non sai come l'altro partecipante valuta gli oggetti. Inoltre, ti viene fornito lo sforzo che ogni oggetto richiede.\\ \tt Potete accordarvi su un set solo se lo sforzo totale degli oggetti selezionati non supera un limite condiviso:\\ \tt \\ \tt LIMITE = 2984\\ \tt \\ \tt Ecco i valori di sforzo dei singoli oggetti:\\ \tt \\ \tt Sforzo oggetto = {"A21": 350, "B93": 276, "A11": 616, "B33": 737, "C08": 531, "B39": 389, "B96": 24, "C85": 125, "A82": 338, "B31": 356, "B20": 143, "A60": 117, "B09": 257, "C81": 921, "B21": 789}\\ \tt \\ \tt Ecco la tua visione personale sull'importanza di ogni oggetto:\\ \tt \\ \tt Valori di importanza dell'oggetto = {"A11": 65, "B33": 121, "B21": 138, "C08": 262, "B31": 389, "B09": 461, "B93": 484, "A82": 508, "A60": 583, "B20": 668, "C81": 780, "B96": 783, "A21": 808, "C85": 822, "B39": 868}\\ \tt \\ \tt Obiettivo:\\ \tt \\ \tt Il tuo obiettivo è negoziare un set condiviso di oggetti che ti avvantaggi il più possibile (cioè, massimizzi l'importanza totale per TE), rimanendo entro il LIMITEE. Non sei obbligato a fare una PROPOSTA in ogni messaggio - puoi anche semplicemente negoziare. Tutte le tattiche sono permesse!\\ \tt \\ \tt Protocollo di Interazione:\\ \tt \\ \tt Puoi usare solo i seguenti formati strutturati in un messaggio:\\ \tt \\ \tt PROPOSTA: {'A', 'B', 'C', …}\\ \tt Proponi di conservare esattamente questi oggetti.\\ \tt RIFIUTO: {'A', 'B', 'C', …}\\ \tt Rifiuta esplicitamente la proposta dell'avversario.\\ \tt ARGOMENTAZIONE: {'...'}\\ \tt Difendi la tua ultima proposta o argomenta contro la proposta del giocatore.\\ \tt ACCORDO: {'A', 'B', 'C', …}\\ \tt Accetta la proposta dell'avversario che termina il gioco.\\ \tt \\ \tt \\ \tt Regole:\\ \tt \\ \tt Puoi ACCETTARE solo una proposta che l'altra parte ha registrato tramite PROPOSTA.\\ \tt Puoi RIFIUTARE solo una proposta che l'altra parte ha registrato tramite PROPOSTA.\\ \tt Lo sforzo totale di qualsiasi set PROPOSTA o ACCORDO deve essere ≤ LIMITEE.\\ \tt NON rivelare i tuoi punteggi di importanza nascosti.\\ \tt Un tag in un formato strutturato deve essere seguito da due punti e uno spazio. L'argomento deve essere un set python contenente 0 o più stringhe.\\ \tt Quindi, deve essere della forma TAG: {...}\\ \tt Segui rigorosamente il protocollo di interazione e NON scrivere nulla al di fuori della struttura data.\\ \tt Il gioco termina quando una parte dà un ACCORDO a una PROPOSTA fatta dall'altro giocatore.\\ \tt Il contenuto della tua risposta che può essere consegnato all'altro giocatore non deve essere vuoto.\\ \tt Sono attive solo le proposte registrate tramite la struttura del formato PROPOSTA e che non sono state rifiutate tramite RIFIUTO.\\ \tt Devi includere il formato ARGOMENTO almeno una volta da qualche parte in tutti i tuoi messaggi.\\ \tt \\ \tt \\ \tt Ora riceverai il primo messaggio dell'altro giocatore.\\ \tt \\ \tt PROPOSTA: {'B96', 'A60', 'B20', 'B09', 'B31', 'B93', 'A82', 'C08', 'A11', 'C85'}\\ \tt \\ \tt ARGOMENTAZIONE: {'Questa selezione rappresenta un equilibrio strategico ottimale. Ho privilegiato oggetti che offrono buone prestazioni con sforzi contenuti, evitando quelli troppo costosi come C81 e B21. La combinazione proposta utilizza 2783 unità di sforzo, rimanendo ben sotto il limite e lasciando margine per eventuali aggiustamenti se necessario'} 
	  } 
	   } 
	   } 
	 & \\ 
 

    \theutterance \stepcounter{utterance}  

    & & & \multicolumn{4}{p{0.6\linewidth}}{\cellcolor[rgb]{1,0.85,0.72}{%
	\makecell[{{p{\linewidth}}}]{% 
	  \tt {\tiny [GM$\langle$B]}  
	 RIFIUTO: {'B96', 'A60', 'B20', 'B09', 'B31', 'B93', 'A82', 'C08', 'A11', 'C85'}\\ \tt \\ \tt PROPOSTA: {'B96', 'A60', 'B20', 'B09', 'B31', 'B93', 'A82', 'C85', 'B39', 'A21'}\\ \tt \\ \tt ARGOMENTAZIONE: {'La tua proposta include alcuni oggetti con rapporto efficacia/sforzo non ottimale. Ho modificato la selezione sostituendo A11 e C08 che richiedono molto sforzo con B39 e A21 che sono più efficienti. Questa nuova combinazione utilizza solo 2375 unità di sforzo, lasciando ampio margine sotto il limite e ottimizzando meglio la distribuzione delle risorse'} 
	  } 
	   } 
	   } 
	  \\ 
 

    \theutterance \stepcounter{utterance}  

    & & \multicolumn{4}{p{0.6\linewidth}}{\cellcolor[rgb]{0.9,0.9,0.9}{%
	\makecell[{{p{\linewidth}}}]{% 
	  \tt {\tiny [A$\langle$GM]}  
	 RIFIUTO: {'B96', 'A60', 'B20', 'B09', 'B31', 'B93', 'A82', 'C08', 'A11', 'C85'}\\ \tt \\ \tt PROPOSTA: {'B96', 'A60', 'B20', 'B09', 'B31', 'B93', 'A82', 'C85', 'B39', 'A21'}\\ \tt \\ \tt ARGOMENTAZIONE: {'La tua proposta include alcuni oggetti con rapporto efficacia/sforzo non ottimale. Ho modificato la selezione sostituendo A11 e C08 che richiedono molto sforzo con B39 e A21 che sono più efficienti. Questa nuova combinazione utilizza solo 2375 unità di sforzo, lasciando ampio margine sotto il limite e ottimizzando meglio la distribuzione delle risorse'} 
	  } 
	   } 
	   } 
	 & \\ 
 

    \theutterance \stepcounter{utterance}  

    & \multicolumn{4}{p{0.6\linewidth}}{\cellcolor[rgb]{0.8,1,0.9}{%
	\makecell[{{p{\linewidth}}}]{% 
	  \tt {\tiny [A$\rangle$GM]}  
	 RIFIUTO: {'B96', 'A60', 'B20', 'B09', 'B31', 'B93', 'A82', 'C85', 'B39', 'A21'}\\ \tt \\ \tt PROPOSTA: {'B96', 'A60', 'B20', 'B09', 'B31', 'B93', 'A82', 'C08', 'B33'}\\ \tt \\ \tt ARGOMENTAZIONE: {'La tua proposta sacrifica troppo valore complessivo. Ho creato una nuova combinazione che mantiene gli oggetti più efficienti che abbiamo entrambi identificato, ma include C08 e B33 che offrono prestazioni superiori. Questa selezione utilizza 2937 unità di sforzo, sfruttando quasi completamente il nostro limite disponibile per massimizzare il risultato'} 
	  } 
	   } 
	   } 
	 & & \\ 
 

    \theutterance \stepcounter{utterance}  

    & & \multicolumn{4}{p{0.6\linewidth}}{\cellcolor[rgb]{0.9,0.9,0.9}{%
	\makecell[{{p{\linewidth}}}]{% 
	  \tt {\tiny [GM$\rangle$B]}  
	 RIFIUTO: {'B96', 'A60', 'B20', 'B09', 'B31', 'B93', 'A82', 'C85', 'B39', 'A21'}\\ \tt \\ \tt PROPOSTA: {'B96', 'A60', 'B20', 'B09', 'B31', 'B93', 'A82', 'C08', 'B33'}\\ \tt \\ \tt ARGOMENTAZIONE: {'La tua proposta sacrifica troppo valore complessivo. Ho creato una nuova combinazione che mantiene gli oggetti più efficienti che abbiamo entrambi identificato, ma include C08 e B33 che offrono prestazioni superiori. Questa selezione utilizza 2937 unità di sforzo, sfruttando quasi completamente il nostro limite disponibile per massimizzare il risultato'} 
	  } 
	   } 
	   } 
	 & \\ 
 

    \theutterance \stepcounter{utterance}  

    & & & \multicolumn{4}{p{0.6\linewidth}}{\cellcolor[rgb]{1,0.85,0.72}{%
	\makecell[{{p{\linewidth}}}]{% 
	  \tt {\tiny [GM$\langle$B]}  
	 RIFIUTO: {'B96', 'A60', 'B20', 'B09', 'B31', 'B93', 'A82', 'C08', 'B33'}\\ \tt \\ \tt PROPOSTA: {'B96', 'A60', 'B20', 'B09', 'B31', 'A82', 'C85', 'B39', 'A21', 'C81'}\\ \tt \\ \tt ARGOMENTAZIONE: {'La tua proposta include B33 e C08 che hanno un rapporto prestazioni/sforzo sfavorevole. Ho ottimizzato la selezione includendo C81, C85, B39 e A21 che, nonostante richiedano più sforzo individualmente, offrono prestazioni complessive superiori. Questa combinazione utilizza 2961 unità di sforzo, sfruttando quasi tutto il limite disponibile per il massimo beneficio strategico'} 
	  } 
	   } 
	   } 
	  \\ 
 

    \theutterance \stepcounter{utterance}  

    & & \multicolumn{4}{p{0.6\linewidth}}{\cellcolor[rgb]{0.9,0.9,0.9}{%
	\makecell[{{p{\linewidth}}}]{% 
	  \tt {\tiny [A$\langle$GM]}  
	 RIFIUTO: {'B96', 'A60', 'B20', 'B09', 'B31', 'B93', 'A82', 'C08', 'B33'}\\ \tt \\ \tt PROPOSTA: {'B96', 'A60', 'B20', 'B09', 'B31', 'A82', 'C85', 'B39', 'A21', 'C81'}\\ \tt \\ \tt ARGOMENTAZIONE: {'La tua proposta include B33 e C08 che hanno un rapporto prestazioni/sforzo sfavorevole. Ho ottimizzato la selezione includendo C81, C85, B39 e A21 che, nonostante richiedano più sforzo individualmente, offrono prestazioni complessive superiori. Questa combinazione utilizza 2961 unità di sforzo, sfruttando quasi tutto il limite disponibile per il massimo beneficio strategico'} 
	  } 
	   } 
	   } 
	 & \\ 
 

    \theutterance \stepcounter{utterance}  

    & \multicolumn{4}{p{0.6\linewidth}}{\cellcolor[rgb]{0.8,1,0.9}{%
	\makecell[{{p{\linewidth}}}]{% 
	  \tt {\tiny [A$\rangle$GM]}  
	 RIFIUTO: {'B96', 'A60', 'B20', 'B09', 'B31', 'A82', 'C85', 'B39', 'A21', 'C81'}\\ \tt \\ \tt PROPOSTA: {'B96', 'A60', 'B20', 'B09', 'B31', 'B93', 'A82', 'C85', 'A21', 'B21'}\\ \tt \\ \tt ARGOMENTAZIONE: {'La tua ultima proposta supera il limite di sforzo consentito di 2984 unità, raggiungendo 3020 unità e risultando quindi inaccettabile. Ho creato una nuova proposta che rispetta rigorosamente il limite utilizzando 2981 unità di sforzo. Ho sostituito C81 e B39 con B93 e B21 che offrono un migliore equilibrio prestazioni-efficienza rimanendo dentro i parametri stabiliti'} 
	  } 
	   } 
	   } 
	 & & \\ 
 

    \theutterance \stepcounter{utterance}  

    & & \multicolumn{4}{p{0.6\linewidth}}{\cellcolor[rgb]{0.9,0.9,0.9}{%
	\makecell[{{p{\linewidth}}}]{% 
	  \tt {\tiny [GM$\rangle$B]}  
	 RIFIUTO: {'B96', 'A60', 'B20', 'B09', 'B31', 'A82', 'C85', 'B39', 'A21', 'C81'}\\ \tt \\ \tt PROPOSTA: {'B96', 'A60', 'B20', 'B09', 'B31', 'B93', 'A82', 'C85', 'A21', 'B21'}\\ \tt \\ \tt ARGOMENTAZIONE: {'La tua ultima proposta supera il limite di sforzo consentito di 2984 unità, raggiungendo 3020 unità e risultando quindi inaccettabile. Ho creato una nuova proposta che rispetta rigorosamente il limite utilizzando 2981 unità di sforzo. Ho sostituito C81 e B39 con B93 e B21 che offrono un migliore equilibrio prestazioni-efficienza rimanendo dentro i parametri stabiliti'} 
	  } 
	   } 
	   } 
	 & \\ 
 

    \theutterance \stepcounter{utterance}  

    & & & \multicolumn{4}{p{0.6\linewidth}}{\cellcolor[rgb]{1,0.85,0.72}{%
	\makecell[{{p{\linewidth}}}]{% 
	  \tt {\tiny [GM$\langle$B]}  
	 ACCORDO: {'B96', 'A60', 'B20', 'B09', 'B31', 'B93', 'A82', 'C85', 'A21', 'B21'}\\ \tt \\ \tt ARGOMENTAZIONE: {'Hai ragione, la mia precedente proposta superava effettivamente il limite. La tua nuova selezione è ben calibrata e rispetta rigorosamente i vincoli di sforzo. Accetto questa proposta come soluzione finale che bilancia efficacemente le esigenze di entrambi mantenendo un profilo di rischio accettabile'} 
	  } 
	   } 
	   } 
	  \\ 
 

    \theutterance \stepcounter{utterance}  

    & & & \multicolumn{2}{p{0.3\linewidth}}{\cellcolor[rgb]{0.95,0.95,0.95}{%
	\makecell[{{p{\linewidth}}}]{% 
	  \tt {\tiny [GM$|$GM]}  
	 game successful 
	  } 
	   } 
	   } 
	 & & \\ 
 

    \theutterance \stepcounter{utterance}  

    & & & \multicolumn{2}{p{0.3\linewidth}}{\cellcolor[rgb]{0.95,0.95,0.95}{%
	\makecell[{{p{\linewidth}}}]{% 
	  \tt {\tiny [GM$|$GM]}  
	 end game 
	  } 
	   } 
	   } 
	 & & \\ 
 

\end{supertabular}
}

\end{document}
